\documentclass[11pt,a4paper]{article}
\usepackage[utf8]{inputenc}
\usepackage[T1]{fontenc}
\usepackage{lmodern}
\usepackage[french]{babel}
\usepackage{amsmath,amssymb,mathtools}
\usepackage{icomma}
\usepackage{xcolor}
\usepackage{tikz}
\usepackage[most]{tcolorbox}
\usepackage{enumitem}
\usepackage{array,colortbl}
\usepackage[a4paper,margin=2cm,headheight=26pt,headsep=10pt]{geometry}
\usepackage{fancyhdr}
\usepackage[hidelinks]{hyperref}
\definecolor{primary}{RGB}{222,49,99}
\definecolor{primarydark}{RGB}{150,22,55}
\definecolor{methcol}{RGB}{52,92,148}
\definecolor{excol}{RGB}{0,135,90}
\definecolor{defcol}{RGB}{0,114,178}
\definecolor{attncol}{RGB}{200,40,55}
\tcbset{boxbase/.style={enhanced,breakable,boxrule=0.9pt,arc=2.5pt,
  left=3mm,right=3mm,top=2.2mm,bottom=2.2mm,fonttitle=\bfseries,coltitle=white}}
\newtcolorbox{cle}[1][]{boxbase,colback=defcol!6,colframe=defcol,
  title={\textsc{À retenir}\ifx\relax#1\relax\relax\else\textmd{ : \itshape #1}\fi}}
\newtcolorbox{methode}[1][]{boxbase,colback=methcol!7,colframe=methcol,sharp corners,
  title={\textsc{Méthode}\ifx\relax#1\relax\relax\else\textmd{ --- \itshape #1}\fi}}
\newtcolorbox{exemplebox}[1][]{boxbase,colback=excol!5,colframe=excol,
  title={\textsc{Exemple}\ifx\relax#1\relax\relax\else\textmd{ : \itshape #1}\fi}}
\newtcolorbox{piege}[1][]{boxbase,colback=attncol!6,colframe=attncol,
  title={\textsc{Attention}\ifx\relax#1\relax\relax\else\textmd{ : \itshape #1}\fi}}
\newcommand{\R}{\mathbb{R}}
\newcommand{\ds}{\displaystyle}
\pagestyle{fancy}\fancyhf{}
\fancyhead[L]{\small\textcolor{primarydark}{\textbf{Thème 3} — Factorisation et résolution}}
\fancyhead[R]{\small\textcolor{primarydark}{\textsc{Tutoriel}}}
\fancyfoot[C]{\small\thepage}
\renewcommand{\headrulewidth}{0.8pt}
\begin{document}

\begin{center}
{\LARGE\bfseries\color{primarydark} Thème 3 — Factorisation et résolution}\\[2pt]
{\color{primary}\rule{0.5\textwidth}{1.2pt}}\\[4pt]
{\small\itshape Tutoriel — la mécanique à maîtriser avant les exercices}
\end{center}

\vspace{1mm}
\noindent \textbf{Factoriser}, c'est récrire une somme en \emph{produit}. Le but : un produit est nul
si et seulement si un facteur est nul, ce qui casse une équation compliquée en équations simples.

\begin{cle}[la règle du produit nul, moteur de tout]
\[ F\times G=0 \iff F=0\ \text{ou}\ G=0. \]
C'est pourquoi on cherche \emph{toujours} à factoriser avant de résoudre une équation.
\end{cle}

\begin{cle}[les trois outils de factorisation]
\begin{enumerate}[leftmargin=1.5em,topsep=1pt,itemsep=1pt]
  \item \textbf{Mise en évidence} : sortir un facteur commun. $\ 2x^{2}-3x=x(2x-3)$.
  \item \textbf{Produits remarquables} :
  \[ A^{2}-B^{2}=(A-B)(A+B);\quad A^{2}\pm 2AB+B^{2}=(A\pm B)^{2}. \]
  \item \textbf{Trinôme du second degré} : si $\Delta=B^{2}-4AC>0$,
  \[ Ax^{2}+Bx+C=A(x-x_1)(x-x_2),\quad x_{1,2}=\frac{-B\pm\sqrt{\Delta}}{2A}. \]
  Si $\Delta=0$ : racine double $x_1=\tfrac{-B}{2A}$. Si $\Delta<0$ : pas de racine réelle.
\end{enumerate}
\end{cle}

\begin{piege}[le discriminant décide sans tout calculer]
Le \emph{signe} de $\Delta=B^{2}-4AC$ donne le nombre de racines réelles : $\Delta>0\Rightarrow$ deux,
$\Delta=0\Rightarrow$ une (double), $\Delta<0\Rightarrow$ aucune. Très utile pour les questions « combien
de solutions~? » sans résoudre.
\end{piege}

\begin{methode}[factoriser un polynôme de degré $\geq 3$ : schéma de Horner]
\begin{enumerate}[leftmargin=1.6em,topsep=1pt,itemsep=0.5pt]
  \item aligner les coefficients (mettre $0$ pour une puissance manquante) ;
  \item trouver une racine évidente $a$ (souvent $\pm1,\pm2$) telle que $P(a)=0$ ;
  \item abaisser le 1\textsuperscript{er} coefficient, puis « $\times a$ + coefficient suivant » de proche en proche ;
  \item les sommes obtenues (sauf la dernière) sont les coefficients de $Q(x)$ ; la dernière est le \emph{reste} (=$0$).
\end{enumerate}
\end{methode}

\begin{exemplebox}[Horner sur $P(x)=x^{3}-x^{2}+2x-2$, racine $a=1$]
\begin{center}\renewcommand{\arraystretch}{1.25}
\begin{tabular}{c|cccc}
 & $1$ & $-1$ & $2$ & $-2$\\
$a=1$ & & $1$ & $0$ & $2$\\ \hline
 & $1$ & $0$ & $2$ & $\mathbf{0}$ (reste)
\end{tabular}
\end{center}
On lit $Q(x)=x^{2}+2$, donc $P(x)=(x-1)(x^{2}+2)$.
\end{exemplebox}

\begin{methode}[simplifier une fraction, résoudre une inéquation]
\textbf{Fraction} $\tfrac{P(x)}{Q(x)}$ : (1) C.E. — annuler $Q$ et l'exclure ; (2) factoriser $P$ et $Q$ ;
(3) simplifier les facteurs communs.\\
\textbf{Inéquation} produit/quotient : factoriser en facteurs du 1\textsuperscript{er} degré, puis dresser
un \textbf{tableau de signes} (chaque facteur $x-r$ est $-$ avant $r$, $+$ après ; on multiplie les
signes colonne par colonne).
\end{methode}

\begin{exemplebox}[tableau de signes de $\tfrac{(x-1)(2x-3)}{x+1}$]
Valeurs critiques $-1$ (interdite), $1$, $\tfrac32$.
\begin{center}\renewcommand{\arraystretch}{1.2}\small
\begin{tabular}{c|ccccc}
$x$ & $-\infty$ & $-1$ & $1$ & $\tfrac32$ & $+\infty$\\ \hline
$\tfrac{(x-1)(2x-3)}{x+1}$ & $-$ & $\nexists$ $+$ & $0$ $-$ & $0$ $+$ &
\end{tabular}
\end{center}
Donc $\tfrac{(x-1)(2x-3)}{x+1}>0$ sur $S=\,]-1,1[\cup\,]\tfrac32,+\infty[$.
\end{exemplebox}

\end{document}
